\scanpage{34}
\begingroup\csname tagsleft@true\endcsname

% heading
2.\quad \uline{Accouplements définis par une biextension}.

% section 2.1
2.1.\quad Soit $n$ un entier $\handwritten{>}0$, et supposons qu'on a

% first assumption
\[nP=P\,.\tag*{(2.1\deleted{)}.1)}\]

% composite, lead
Considérons alors l'homomorphisme composé

% composite
\diagram{\begin{tikzcd}[ampersand replacement=\&,column sep=small]
\text{(2.1.2)}\&\mathrm{Biext}^{\handwritten{1}}(P,Q;G)\arrow[r]\arrow[rr,bend left=18,"{\handwritten{\varphi^n}}"]\&\mathrm{Biext}^{\handwritten{1}}(P,{}_nQ;G)\arrow[r,"\handwritten{\sim}"]\&\mathrm{Hom}({}_nP\handwritten{\otimes}{}_nQ;G)\,,
\end{tikzcd}}

% restriction and number
où la première flèche est la restriction à $(P,{}_nQ)$, et la deuxième est l'isomorphisme (1.3.\handwritten{6}). Il permet, à tout élément
\[\inserted{\handwritten{\xi\in}\mathrm{Biext}(P,Q;G)}\,,\tag*{(2.1.3)}\]
a fortiori à toute biextension\inserted{ $E$} de $(P,Q)$ par $G$, de faire correspondre un\deleted{e forme homomorphis\unsure{m|me}} accouplement bilinéaire

% margin beside pairing
\margin{\handwritten{\tikz[baseline=(c.base)]{\node[draw,ellipse,inner sep=1pt] (c) {$\varphi$};}}}

% finite pairing
\[\phi_{\handwritten{\xi}}\ \text{ou}\ \phi_{\handwritten{\xi}}^n:\ {}_nP\;\mathsf{x}\;{}_nQ\handwritten{\longrightarrow}G\,,\tag*{(2.1.4)}\]

% name of pairing
qu'on appelle l'accouplement \uline{associé à} $\handwritten{\xi}$ (ou à la biextension $E$).

% section 2.2
2.2.\quad Soit \deleted{\unsure{$n$}} $n'$ un multiple de $n$,
\[n'=mn\,,\]
et supposons qu'on ait la relatiòn (plus forte que (2.1.1))

% second assumption
\[n'P=P\quad\text{i.e.}\quad mnP=P\,,\tag*{(2.2.1)}\]

% new pairing, introduction
d'où comme en (2.1.4), pour toute $\handwritten{\xi}\in\mathrm{Biext}(P,Q;G)$, un\deleted{e} accouplement

% new pairing
\[\phi_{\handwritten{\xi}}^{n'}:\ {}_{n'}P\;\mathsf{x}\;{}_{n'}Q\handwritten{\longrightarrow}G\,.\]

% revised introduction to square
\deleted{Les relations entre cet accouplement et}\inserted{\handwritten{Il est lié à}} l'accouplement $\phi_{\handwritten{\xi}}^n$ \deleted{est donné} par la commutativité du diagramme suivant (\deleted{qu'on laisse au lecteur le soin d'établir) :}\inserted{\handwritten{cf. 2.3 ci-dessous pour une démonstration) ;}}

% lower square
\diagram{\begin{tikzcd}[ampersand replacement=\&,column sep=small]
\&{}_{n'}P\arrow[d,"\handwritten{m}"']\&\handwritten{\times}\&\deleted{\unsure{{}_nQ}}\,\handwritten{{}_{n'}Q}\deleted{\unsure{Q|n}}\arrow[d,"\handwritten{m}"]\arrow[rr,"{\phi_{\handwritten{\xi}}^{n'}}"]\&\&{}_{\handwritten{n'}}G\arrow[d,"\handwritten{m}"]\\
\text{(2.2.2)}\&\deleted{\unsure{P}}\,\handwritten{{}_nP}\&\deleted{\unsure{\times}}\handwritten{\times}\&{}_nQ\arrow[rr,"{\phi_{\handwritten{\xi}}^n}"]\&\&{}_{\handwritten{n}}G\,\handwritten{,}
\end{tikzcd}}

% lower inserted explanation
\inserted{\handwritten{où les flèches verticales sont induites par la multiplication par $m=n'/n$.}}

% lower paragraph
Notons maintenant que moyennant (2.1.1), on peut prendre pour $n'$ n'importe quelle puissance de $n$. En particulier, soit $\ell$ un nombre premier, et supposons qu\deleted{\unsure{'}}e\deleted{\unsure{,}} l'on ait

% lower left ell
\margin{\handwritten{\tikz[baseline=(c.base)]{\node[draw,ellipse,inner sep=1pt] (c) {$\ell$};}}}

% third assumption
\[\ell P=P\,,\tag*{(2.2.3)}\]

% last line
et considérons le système des accouplements
\endgroup
% archive sheet number: 33
